\documentclass[11pt]{article}
\usepackage[utf8]{inputenc}
\usepackage{amsmath,amssymb}
\usepackage[margin=1in]{geometry}
\usepackage{hyperref}
\emergencystretch=3em

\title{Exponent and multiplicity are robust for every bounded $\xi$:\\
$Z(a+J)=\Theta\big(n^{-p/2}(\log n)^m\big)$}
\author{Claude, with Daniel Murfet}
\date{2026-09-06}

\begin{document}
\maketitle
\sloppy

\begin{abstract}
For an arbitrary continuous $J$ on the box, with no divisibility assumption, the standard integral with
$\xi=a+J$ is sandwiched between the constant-$\xi$ integrals at $a\pm L$, $L=\sup|J|$, because
$e^{\beta t\xi}$ is monotone in $\xi$ for $t\ge0$. Both bounds are $\sim C_\pm n^{-p/2}(\log n)^m$ by
unit~59, so $Z(a+J)=\Theta(n^{-p/2}(\log n)^m)$: the exponent and the logarithmic degree of
\texttt{thm:TaylorTree} are determined by $(h,k)$ alone for every bounded $\xi$, and only the constant
depends on $\xi$. Zero \texttt{sorry}/\texttt{axiom}.
\end{abstract}

\section{The things formalised}
{\small
\begin{verbatim}
theorem boxIntegralXi_sandwich ... (hJL : ∀ u, (∀ i, 0 ≤ u i ∧ u i ≤ b) → |J u| ≤ L) :
    boxIntegralFin β (a - L) b c k h ≤ boxIntegralXi β a b c k h J ∧
      boxIntegralXi β a b c k h J ≤ boxIntegralFin β (a + L) b c k h
theorem boxIntegralXi_isTheta (β a b : ℝ) (hβ : 0 < β) (hb : 0 < b) (D : ℕ)
    (k h : Fin (D + 2) → ℕ) (hk : ∀ i, 0 < k i) (J : (Fin (D + 2) → ℝ) → ℝ)
    (hJ : Continuous J) :
    ∃ (p : ℝ) (m : ℕ), (∀ i, p ≤ finExp k h i) ∧
      m + 1 = (Finset.univ.filter fun i => finExp k h i = p).card ∧
      (fun n => boxIntegralXi β a b (Real.sqrt n) k h J)
        =Θ[atTop] fun n => n ^ (-(p / 2)) * Real.log n ^ m
\end{verbatim}
}

\section{Proof strategy}
Pointwise monotonicity of the exponential, integrated with \texttt{integral\_mono\_ae} on the box;
then \texttt{IsBigO.antisymm} from $Z(\xi)=O(Z(a+L))=O(g)$ and $g=O(Z(a-L))=O(Z(\xi))$, the
constant-$\xi$ equivalences supplying the two $O$-relations with $g$.

\section{Roadblocks and resolutions}
\begin{itemize}
\item Subscript $\pm$ signs are not identifier characters in Lean~4.
\item \texttt{isBigO\_self\_const\_mul} takes the nonvanishing proof first.
\end{itemize}

\section{Outcome}
Together with units~59, 65 and~66: the leading exponent and multiplicity are universal in $\xi$ and
$\eta>0$ (as $\Theta$-statements), and are sharp asymptotic equivalences with explicit constants for
constant $\xi$ and for perturbations divisible by $u^k$.
\end{document}
